\documentclass[letterpaper]{article}
\usepackage{amssymb}
\usepackage{fullpage}
\usepackage{amsmath}
\usepackage{epsfig,float,alltt}
\usepackage{psfrag,xr}
\usepackage[T1]{fontenc}
\usepackage{url}
\usepackage{pdfpages}
\usepackage{halloweenmath}
%\includepdfset{pagecommand=\thispagestyle{fancy}}
\author{Prof. Grizzle}
\title{Rob 501 - Mathematics for R$\skull$b$\skull$tics\\
HW \#7}

\begin{document}
\date{}
\maketitle


 \newcommand{\trace}{\mathrm{trace}}
\newcommand{\real}{\mathbb R}  % real numbers  {I\!\!R}
\newcommand{\nat}{\mathbb N}   % Natural numbers {I\!\!N}
\newcommand{\cp}{\mathbb C}    % complex numbers  {I\!\!\!\!C}
\newcommand{\ds}{\displaystyle}
\newcommand{\mf}[2]{\frac{\ds #1}{\ds #2}}
\newcommand{\Nagy}[2]{{Nagy, Page~#1, }{Prob.~#2}}
\newcommand{\spanof}[1]{\textrm{span} \{ #1 \}}
\parindent 0pt


\begin{flushright}
{\bfseries Due Nov. 9, 2017} \\
{\bfseries 3PM via Gradescope} \\[3mm]
\end{flushright}

\noindent \textbf{Remarks:}  Problems 1 through 5 involve calculations. Problem 5 is very important. Please spend extra time on it as it is very helpful for understanding the Kalman Filter. Problem 6 explains why we can often obtain recursion relations of the type: $ \widehat{x}_{k+1}$ is a linear combination of $\widehat{x}_k$ and the ``innovation'' or new measured information $ (y_{k+1} - \widehat{y}_{k+1|k} )$. If you are pressed for time, skip Problem 6 and study the solutions. It is better to spend your time on Problem 5.

\begin{enumerate}


%Prob. 1 Matrices
\item Classify each matrix as positive definite, positive semi-definite, or neither. In addition, if the matrix is either positive definite or positive semi-definite, find a square root. You may use MATLAB to factor a symmetric matrix as $\Lambda = O^\top P O$ or as $P= O \Lambda O ^\top$.

\begin{enumerate}
\setlength{\itemsep}{.1in}
\renewcommand{\labelenumi}{(\alph{enumi})}

\item $\greatpumpkin= \left[ \begin{array}{rr}  1&   3\\  3&   9\end{array} \right].$

\item  $\bigskull= \left[ \begin{array}{rrr}  6&  10&  11\\ 10&  19&  19\\ 11&  19&  21\end{array} \right].$

\item $\mathghost=  \left[ \begin{array}{rrr}  2&   6&  10\\  6&  10&  14\\ 10&  14&  18\end{array} \right].$

\end{enumerate}


%Prob. 2 Matrices
\item Use the results on Schur Complements  to solve  the following problems BY HAND:

        \begin{enumerate}
\setlength{\itemsep}{.1in}
\renewcommand{\labelenumi}{(\alph{enumi})}

\item  Determine if  $\greatpumpkin= \left[ \begin{array}{rr}  1&   3\\  3&   8\end{array} \right]$ is positive definite or not.

\item  Determine if  $\bigskull= \left[ \begin{array}{rrr}  1&   0&  6\\  0&  4&   7\\  6&   7&  10\end{array} \right]$ is positive definite or not.

\item Find the range of $a$ such that the following matrix is positive definite:
$\mathghost= \left[ \begin{array}{rrr}  1&   2&   6\\  2&   5&   7\\  6&   7&  a\end{array}\right]$


\end{enumerate}

%Prob. Underdetermine equations
\item Find $x$ of minimum norm that satisfies the equation

$$ \left[ \begin{array}{rrr}  1&   3 & 2\\  3&   8 & 4 \end{array} \right] x = \left[ \begin{array}{r}  1 \\2\end{array} \right]$$

        \begin{enumerate}
\setlength{\itemsep}{.1in}
\renewcommand{\labelenumi}{(\alph{enumi})}

\item  Use the standard inner product on $\real^3$

\item  Use the inner product $<x,y>= x^\top \left[ \begin{array}{rrr}  5&   1&   9\\  1&   2&   1\\  9&   1&  17\end{array} \right]y$

\end{enumerate}

 %Prob. 4  BLUE
\item $x$ has been bewitched to give the following data:
$$ y = \mathwitch x + \epsilon$$
with
$$ \mathwitch=\left[\begin{array}{rr}  1&   2\\  3&   4\\  5&   0\\  0&   6\end{array} \right] ~~ y=\left[  \begin{array}{r}1.5377\\3.6948\\-7.7193\\7.3621\end{array}\right] ~~\mbox{and}  ~~ E\{ \epsilon \epsilon^\top \} = Q=\left[ \begin{array}{rrrr}1.00& 0.50& 0.50& 0.25\\0.50& 2.00& 0.25& 1.00\\0.50& 0.25& 2.00& 1.00\\0.25& 1.00& 1.00& 4.00\end{array}
\right] $$
As in class, $E\{ \epsilon\} = 0$.


        \begin{enumerate}
\setlength{\itemsep}{.1in}
\renewcommand{\labelenumi}{(\alph{enumi})}

\item Find the Best Linear Unbiased Estimate (BLUE) for $x$, using only the first two values of $y$. Also compute the covariance of the estimate.

\item Find the Best Linear Unbiased Estimate (BLUE) for $x$, using only the first three values of $y$. Also compute the covariance of the estimate.

\item Find the Best Linear Unbiased Estimate (BLUE) for $x$, using all the values of $y$. Also compute the covariance of the estimate.

\end{enumerate}

\noindent \textbf{Note:} For (a), you use the first 2 rows of $y$ and $\mathwitch$ and the upper $2 \times 2 $ part of $Q$. For (b), you use the first 3 rows of $y$ and $C$, as well as the upper $3 \times 3 $ part of $Q$. You see the pattern, I hope. Do all the calculations in MATLAB. You do not have to turn in your code.

% Prob. 5 Jointly Gaussian Random Variables
\item Read the Handout \texttt{GaussianRandomVariablesAndVectors.pdf}, which you can find on CANVAS. We consider three jointly normal random variables $(X,Y,Z)$, with
$$ \mbox{mean}~~\mu = \left[\begin{array}{r} -1\\  0\\  1\end{array} \right] ~~\mbox{and covariance}~~ \Sigma = \left[  \begin{array}{rrr}  2&   2&   1\\  2&   4&   2\\  1&   2&   2\end{array}
\right] $$

        \begin{enumerate}
\setlength{\itemsep}{.1in}
\renewcommand{\labelenumi}{(\alph{enumi})}

\item  Compute the conditional distribution of  $  \left[\begin{array}{r} X\\  Y\end{array} \right]  \left| {\{Z=z\} } \right.  $, the conditional distribution of the vector $[X,Y]^\top$ given $Z=z$, which is the same as the joint distribution of the normal random variables $X|\{Z=z\}$ and $Y|\{Z=z\}$. To be extra clear, give the mean vector and covariance matrix for $[X,Y]^\top$ given $Z=z$.

    \item Compute the distribution of $X|_{\{Z=z\}}$ conditioned on $Y|_{\{ Z=z\} }= y$.

\item  Compute the conditional distribution of  $  X | _{\left[\begin{array}{r} Y=y\\  Z=z\end{array} \right] } $, or more compactly, $ X_{|Y=y,Z=z}$, the conditional distribution of  $X$ given the vector $[Y=y,Z=z]^\top$.

    \item Compare your answers for (b) and (c).
\end{enumerate}

\newpage

 %Prob. 6  Recursion
\item  Let $({\cal X}, \real, <\cdot, \cdot>)$ be a finite-dimensional inner product space. Our objective is to understand recursion relations similar to what we see in RLS and the Kalman filter. With that in mind, let $\{y_1, \cdots, y_N \}$ be a linearly independent set in ${\cal X}$. For $1 \le k \le N$ define     $$M_k: = \spanof{y_1, \cdots, y_k}$$
      and for $x \in {\cal X}$ define $ \widehat{x}_k := \mathrm{arg~min}_{m\in M_k} ||x-m||,$ which is the orthogonal projection of $x$ onto $M_k$.
  \begin{enumerate}
\setlength{\itemsep}{.1in}
\renewcommand{\labelenumi}{(\alph{enumi})}

\item  Suppose that for some $1 \le k < N$, we have $y_{k+1} \perp M_k$ (interpretation: the new measurement $y_{k+1}$ is orthogonal to the previous measurements). Show that there exists $\beta \in \real$ such that
$$\widehat{x}_{k+1} = \widehat{x}_k + \beta y_{k+1},$$
and give a formula for $\beta$.

\item We no longer make any hypothesis about $y_{k+1}$ being orthogonal to $M_k$. What we do now is, for each $1 \le k < N$, define
$$\widehat{y}_{k+1|k} = \mathrm{arg~min}_{m\in M_{k}} ||y_{k+1}-m||, $$
which we interpret as the orthogonal projection of the new measurement $y_{k+1}$ onto the subspace generated by the previous measurements.
Show that there exists $\beta \in \real$ such that
$$\widehat{x}_{k+1} = \widehat{x}_k + \beta (y_{k+1} - \widehat{y}_{k+1|k} ),$$
and give a formula for $\beta$. \\

\textbf{Remark:} The error term $ (y_{k+1} - \widehat{y}_{k+1|k} )$ is the ``innovations'' in the case of the Kalman filter and RLS. What is particularly nice in the case of RLS and the Kalman filter, where we have a model such as $y_k = C_k x_k + v_k$, we can compute $\widehat{y}_{k+1|k}$ directly from $ \widehat{x}_k$, often with a formula such as $\widehat{y}_{k+1|k} = C_{k+1}  \widehat{x}_k$. In other words, we do not have to solve an extra optimization problem in order to compute $\widehat{y}_{k+1|k};$ instead, we can bootstrap from the previous solution to the optimization problem.
\end{enumerate}



%End of list of HW Problems
\end{enumerate}

\newpage

\begin{center}
{\bf Hints}
\end{center}

\noindent \textbf{Hints: Prob. 1 }  Use the \texttt{help eig} command in MATLAB.

\vspace{1cm}

\noindent \textbf{Hints: Prob. 2}  Recall that for a symmetric matrix $M=\left[  \begin{array}{cc} A & B \\
B^\top & C \end{array} \right]$
the following are equivalent:
 \begin{enumerate}
\setlength{\itemsep}{.1in}
\renewcommand{\labelenumi}{(\alph{enumi})}

\item $M \succ 0$

\item $A \succ 0$ and $C-B^\top A^{-1} B \succ 0$

\item $C \succ 0$ and $A-B C^{-1} B^\top \succ 0$

\end{enumerate}

\vspace{1cm}

\noindent \textbf{Hints: Prob. 3 } This is an under determined system of equations and not an over determined system of equations.

\vspace{1cm}

\noindent \textbf{Hints: Prob. 4 } Recall our formulas (using $C := \mathwitch$ for convenience):
$$ \widehat{K}=(C^\top Q^{-1}C)^{-1} C^\top Q^{-1}~~\mbox{and}~~E\{ (\widehat{x}-x)(\widehat{x}-x)^\top \}=(C^\top Q^{-1}C)^{-1}  $$
For (a), you use the first 2 rows of $y$ and $C$ and the upper $2 \times 2 $ part of $Q$. For (b), you use the first 3 rows of $y$ and $C$, as well as the upper $3 \times 3 $ part of $Q$. You see the pattern, I hope. Do all the calculations in MATLAB. You do not have to turn in your code.
\vspace{1cm}


\noindent \textbf{Hints: Prob. 5 } Print out the handout on Jointly Gaussian Random Vectors, and read it carefully.  Note \textbf{Fact 1:} Conditional Distributions of Gaussian Random Vectors

 \begin{enumerate}
\setlength{\itemsep}{.1in}
\renewcommand{\labelenumi}{(\alph{enumi})}

\item Identify $X_1 =  \left[\begin{array}{r} X \\ Y\end{array} \right] $ and $X_2 = Z$. Based on this, identify and write down $\Sigma_{11}$, $\Sigma_{12}$, $\Sigma_{21}$, and $\Sigma_{22}$, and then $\mu_1$ and $\mu_2$, and finally, note that $x_2=z$. Now apply the formulas for $\mu_{1|2}$ and $\Sigma_{1|2}$. These are the mean and covariance of the jointly normally distributed random variables $X_{|Z=z}$ and $Y_{|Z=z}$.

  \item From (a), we know   $X_{|Z=z}$ and $Y_{|Z=z}$ are jointly distributed normal random variables, and we know their mean and covariance. Rename the mean $\mu$ and the covariance $\Sigma$ (i.e., $\mu_{1|2} \to \mu$ and $\Sigma_{1|2} \to \Sigma$). Using Fact 1, identify $X_1 =  X_{|Z=z}$ and $X_2 = Y_{|Z=z}$, and then  identify and write down $\Sigma_{11}$, $\Sigma_{12}$, $\Sigma_{21}$, and $\Sigma_{22}$, and then $\mu_1$ and $\mu_2$, and finally, note that $x_2=y$.  Now apply the formulas for $\mu_{1|2}$ and $\Sigma_{1|2}$. These are the mean and covariance of a normally distributed random variable. Which one? If you do not know, work part (c) and then return here. If you do know, still work part (c).

  \item Go back to the very beginning with our three jointly normal random variables, and this time identify $X_1 =  X$ and $X_2 = \left[\begin{array}{r} Y \\ Z\end{array} \right] $. Based on this, identify and write down $\Sigma_{11}$, $\Sigma_{12}$, $\Sigma_{21}$, and $\Sigma_{22}$, and then $\mu_1$ and $\mu_2$, and finally, note that $x_2=\left[\begin{array}{r} y \\ z\end{array} \right] $. Now apply the formulas for $\mu_{1|2}$ and $\Sigma_{1|2}$. These are the mean and covariance of the normally distributed random variable $X_{|Y=y,Z=z}$.

      \item Read again the handout on Jointly Gaussian Random Vectors. Go back to the beginning and repeat as necessary: \textbf{FACT 4:} \textit{If we have jointly distributed normal random vectors, when we condition one block of vectors on another, we always obtain either a jointly distributed normal random vector or, if only a scalar quantity is left,  a normally distributed random variable.} This is an amazingly useful property of Gaussian (i.e., normal) random variables.

\end{enumerate}

\vspace{1cm}


\noindent \textbf{Hints: Prob. 6 } There are several ways to approach part (a):
\begin{itemize}
\item Method 1: Let $G_k$ be the Gram matrix for $M_k$ and $G_{k+1}$ be the Gram matrix for $M_{k+1}$. Then, using $y_{k+1} \perp M_k$, you deduce that
    $$G_{k+1} = \left[ \begin{array}{cc} G_k & 0_{1 \times k} \\ 0_{k \times 1} & <y_{k+1}, y_{k+1}> \end{array} \right] $$
    Use this block diagonal structure to relate the solution of the normal equations for $M_{k+1}$ to the solution of the normal equations for $M_{k}$.

\item Method 2: The solution of the optimization problem depends on the subspaces $M_k$ and not on the bases you use for them.  Apply Gram Schmidt to produce an ortho\textbf{normal} basis for $M_k$. Relate it to an orthornormal basis for $M_{k+1}$. Then use what you know about the orthogonal projection of $x$ onto sets with orthonormal bases.

\end{itemize}

For part (b), we know from the Projection Theorem that $y_{k+1} - \widehat{y}_{k+1|k}$ is orthogonal to $M_k$. If you need a second hint, note that
$$M_{k+1} = M_k \oplus \spanof{ y_{k+1} }  =  M_k \oplus \spanof{ y_{k+1} - v },$$
for any $v\in M_k$. Hence, because $\widehat{y}_{k+1|k} \in M_k$ by definition, we have
$$M_{k+1} = M_k \oplus \spanof{ y_{k+1} - \widehat{y}_{k+1|k} }~~\mbox{and}~~ M_k \perp (y_{k+1} - \widehat{y}_{k+1|k}).$$
You can now apply your result from (a).

\newpage


\end{document}


